%%%%%%%%%%%%%%%%
%%% deut 5 %%%
%%%%%%%%%%%%%%%%

\Chapter{5}

\Verse{1}
{
	\hDtrA{וַיִּקְרָא מֹשֶׁה אֶל כּל יִשְׂרָאֵל וַיֹּאמֶר אֲלֵהֶם שְׁמַע יִשְׂרָאֵל אֶת הַחֻקִּים וְאֶת הַמִּשְׁפָּטִים אֲשֶׁר אָנֹכִי דֹּבֵר בְּאזְנֵיכֶם הַיּוֹם וּלְמַדְתֶּם אֹתָם וּשְׁמַרְתֶּם לַעֲשֹׂתָם׃}
}
{
	\eDtrA{
		5 And Moses called all of Israel and said to them,
		“Listen, Israel, to the laws and the judgments that I’m
		speaking in your ears today. And you shall learn them and
		be watchful to do them.
	}
}
{
%	\footnote{
%		watchful. This is the fifth time that Moses uses this word in his
%		speech. Its concentration here calls to mind another element of the
%		beginning of the Torah, where forms of this word are played upon: Humans
%		are placed in the garden of Eden “to watch over it” (Gen 2:15); but when
%		they eat from the forbidden tree they are removed from the garden, and
%		cherubs are placed “to watch over the way to the tree of life”
%		(3:24)—that is, to keep humans away; and in the next episode Cain lies
%		about his brother Abel’s whereabouts by saying “Am I my brother’s
%		watchman?” (4:9). The point of this: The word “watch” marks a gradual
%		descent by humans away from their initial condition of closeness to God.
%		But then the word is used in a new way when God says that the merit of
%		Abraham is that he “kept my watch” (26:5). (This phrase is generally
%		translated “keep my charge” in this work, but I used the word “watch” in
%		Gen 26:5 in order to make known this development in Genesis.) The word
%		thus comes to point the way to a return to harmony with the creator, and
%		thus it resurfaces meaningfully here in the Torah’s final book.
%	}
}

\Verse{2}
{
	\hDtrA{יְהֹוָה אֱלֹהֵינוּ כָּרַת עִמָּנוּ בְּרִית בְּחֹרֵב׃}
}
{
	\eDtrA{YHWH, our {\God}, had made a covenant with us at Horeb.}
}
{
}

\Verse{3}
{
	\hDtrA{לֹא אֶת אֲבֹתֵינוּ כָּרַת יְהֹוָה אֶת הַבְּרִית הַזֹּאת כִּי אִתָּנוּ אֲנַחְנוּ אֵלֶּה פֹה הַיּוֹם כֻּלָּנוּ חַיִּים׃}
}
{
	\eDtrA{
		YHWH did not make this covenant with our fathers, but with
		us! We! These! Here! Today! All of us! Living!
	}
}
{
%	\footnote{
%		us! We! These! Here! Today! All of us! Living! As I commented above
%		(1:9), Moses mixes past, present, and future. He speaks to the people in
%		front of him as if they had all been at Sinai forty years earlier. Now
%		he says it explicitly, powerfully, unmistakably, with seven different
%		words: Each generation must see themselves as personally standing at
%		Sinai, not just as inheriting their parents’ covenant, but as making the
%		covenant themselves. It is a present, living commitment.
%	}
}

\Verse{4}
{
	\hDtrA{פָּנִים בְּפָנִים דִּבֶּר יְהֹוָה עִמָּכֶם בָּהָר מִתּוֹךְ הָאֵשׁ׃}
}
{
	\eDtrA{
		Face-to-face, YHWH spoke with you at the mountain from
		inside the fire
	}
}
{
%	\footnote{
%		Face-to-face. Just a short time earlier Moses had said, “You didn’t see
%		any form in the day that YHWH spoke to you at Horeb from inside the
%		fire.” But now he says, “Face-to-face, YHWH spoke with you at the
%		mountain from inside the fire.” This is further confirmation that
%		“face-to-face” is not a literal expression. Like many Hebrew expressions
%		involving the word “face,” it is an idiom. It is used metaphorically to
%		mean direct, personal communication, not mediated by any third party.
%	}
}

\Verse{5}
{
	\hDtrA{אָנֹכִי עֹמֵד בֵּין יְהֹוָה וּבֵינֵיכֶם בָּעֵת הַהִוא לְהַגִּיד לָכֶם אֶת דְּבַר יְהֹוָה כִּי יְרֵאתֶם מִפְּנֵי הָאֵשׁ וְלֹא עֲלִיתֶם בָּהָר לֵאמֹר׃}
}
{
	\eDtrA{
		(I was standing between YHWH and you at that time to tell
		you YHWH’s word because you were afraid on account of the
		fire, and you didn’t go up in the mountain), saying:
	}
}
{
}

\Verse{6}
{
	\hDtrA{אָנֹכִי יְהֹוָה אֱלֹהֶיךָ אֲשֶׁר הוֹצֵאתִיךָ מֵאֶרֶץ מִצְרַיִם מִבֵּית עֲבָדִ͏ים׃}
}
{
	\eDtrA{
		I am YHWH, your {\God}, who brought you out from the land of
		Egypt, from a house of slaves.
	}
}
{
}

\Verse{7}
{
	\hDtrA{לֹא יִהְיֶ͏ה לְךָ אֱלֹהִים אֲחֵרִים עַל פָּנָ͏ַי׃}
}
{
	\eDtrA{You shall not have other gods before my face.}
}
{
%	\footnote{
%		not have other gods. This is a monotheistic command. It has been
%		mistakenly taken to mean that other gods do exist but that Israelites
%		are not to worship them. But that is a linguistic misunderstanding. See
%		the comment on Exod 20:3.
%	}
}

\Verse{8}
{
	\hDtrA{לֹא תַעֲשֶׂה לְךָ פֶסֶל כּל תְּמוּנָה אֲשֶׁר בַּשָּׁמַיִם מִמַּעַל וַאֲשֶׁר בָּאָרֶץ מִתָּ͏ַחַת וַאֲשֶׁר בַּמַּיִם מִתַּחַת לָאָרֶץ׃}
}
{
	\eDtrA{
		You shall not make a statue, any form that is in the skies
		above or that is in the earth below or that is in the
		water below the earth.
	}
}
{
}

\Verse{9}
{
	\hDtrA{לֹא תִשְׁתַּחֲוֶה לָהֶם וְלֹא תעבְדֵם כִּי אָנֹכִי יְהֹוָה אֱלֹהֶיךָ אֵל קַנָּא פֹּקֵד עֲוֺן אָבוֹת עַל בָּנִים וְעַל שִׁלֵּשִׁים וְעַל רִבֵּעִים לְשֹׂנְאָי׃}
}
{
	\eDtrA{
		You shall not bow to them, and you shall not serve them.
		Because I, YHWH, your {\God}, am a jealous {\God}, counting
		parents’ crime on children and on the third generation and
		on the fourth generation of those who hate me,
	}
}
{
}

\Verse{10}
{
	\hDtrA{וְעֹשֶׂה חֶסֶד לַאֲלָפִים לְאֹהֲבַי וּלְשֹׁמְרֵי (מצותו) [מִצְוֺתָי]׃}
}
{
	\eDtrA{
		but practicing kindness to thousands for those who love me
		and for those who observe my commandments.
	}
}
{
}

\Verse{11}
{
	\hDtrA{לֹא תִשָּׂא אֶת שֵׁם יְהֹוָה אֱלֹהֶיךָ לַשָּׁוְא כִּי לֹא יְנַקֶּה יְהֹוָה אֵת אֲשֶׁר יִשָּׂא אֶת שְׁמוֹ לַשָּׁוְא׃}
}
{
	\eDtrA{
		You shall not bring up the name of YHWH, your {\God}, for a
		falsehood, because YHWH will not make one innocent who
		will bring up His name for a falsehood.
	}
}
{
}

\Verse{12}
{
	\hDtrA{שָׁמוֹר אֶת יוֹם הַשַּׁבָּת לְקַדְּשׁוֹ כַּאֲשֶׁר צִוְּךָ יְהֹוָה אֱלֹהֶיךָ׃}
}
{
	\eDtrA{
		Observe the Sabbath day, to make it holy, as YHWH, your
		{\God}, commanded you.
	}
}
{
%	\footnote{
%		Observe. Probably the most famous difference between the wording of the
%		Ten Commandments at Sinai and the way that Moses quotes them here is
%		this first word of the Sabbath commandment. At Sinai God says “Remember
%		the Sabbath” (Exod 20:8), but Moses now says “Observe the Sabbath.” Many
%		have commented on the two words. The most basic explanation may be that
%		the two are both necessary and are complementary: In the mind, one must
%		remember it. In actions, one must observe it. Another explanation: The
%		most striking difference between the Decalogue at Sinai and Moses’
%		quotation of it here is the reason for the Sabbath command. At Sinai it
%		is Because for six days YHWH made the skies and the earth, the sea, and
%		everything that is in them, and He rested on the seventh day. On account
%		of this, YHWH blessed the Sabbath day and made it holy. (Exod 20:11) But
%		here it is so you shall remember that you were a slave in the land of
%		Egypt and YHWH, your God, brought you out from there with a strong hand
%		and an outstretched arm. On account of this, YHWH, your God, has
%		commanded you to do the Sabbath day. (Deut 5:15) What do we learn from
%		this: When God states the commandment, God identifies the Sabbath with
%		the creation of the universe and the sanctification of time (see the
%		comments on Gen 2:2 and 2:3). But when a human, Moses, states the
%		commandment, he identifies the Sabbath rather with history, with a
%		divine action in human affairs. The Sabbath thus comes to have both
%		dimensions: cosmic and historical. (A Dead Sea scroll text of
%		Deuteronomy combines these and gives both reasons for the Sabbath.) This
%		in turn may explain why the text here begins this command with the words
%		“Observe the Sabbath” instead of “Remember the Sabbath”—namely, since it
%		now adds the words “and you shall remember that you were a slave …” at
%		the end of the commandment, it uses a different word at the beginning.
%	}
}

\Verse{13}
{
	\hDtrA{שֵׁשֶׁת יָמִים תַּעֲבֹד וְעָשִׂיתָ כּל מְלַאכְתֶּךָ׃}
}
{
	\eDtrA{Six days you shall labor and do all your work,}
}
{
}

\Verse{14}
{
	\hDtrA{וְיוֹם הַשְּׁבִיעִי שַׁבָּת לַיהֹוָה אֱלֹהֶיךָ לֹא תַעֲשֶׂה כל מְלָאכָה אַתָּה וּבִנְךָ וּבִתֶּךָ וְעַבְדְּךָ וַאֲמָתֶךָ וְשׁוֹרְךָ וַחֲמֹרְךָ וְכל בְּהֶמְתֶּךָ וְגֵרְךָ אֲשֶׁר בִּשְׁעָרֶיךָ לְמַעַן יָנוּחַ עַבְדְּךָ וַאֲמָתְךָ כָּמוֹךָ׃}
}
{
	\eDtrA{
		and the seventh day is a Sabbath to YHWH, your {\God}. You
		shall not do any work: you and your son and your daughter
		and your servant and your maid and your ox and your ass or
		any animal and your alien who is in your gates—in order
		that your servant and your maid will rest like you,
	}
}
{
%	\footnote{
%		you and your son and your daughter. The text does not mention “your
%		wife.” One might take this to be male-chauvinist, but I think that, in
%		this case, it is the opposite. The text mentions daughters as well as
%		sons, and it mentions female as well as male servants here (and
%		elsewhere: Deut 12:12,18). Therefore the words “you and your son and
%		your daughter” must be understood to be directed to both men and women,
%		husbands and wives, in the first place. Indeed, all of the Ten
%		Commandments are directed to each individual Israelite, male and female,
%		with the exception of the first part of the last commandment—“You shall
%		not covet your neighbor’s wife”—which is explicitly directed only to men
%		and reflects a marital structure that men control.
%	}
}

\Verse{15}
{
	\hDtrA{וְזָכַרְתָּ כִּי עֶבֶד הָיִיתָ בְּאֶרֶץ מִצְרַיִם וַיֹּצִאֲךָ יְהֹוָה אֱלֹהֶיךָ מִשָּׁם בְּיָד חֲזָקָה וּבִזְרֹעַ נְטוּיָה עַל כֵּן צִוְּךָ יְהֹוָה אֱלֹהֶיךָ לַעֲשׂוֹת אֶת יוֹם הַשַּׁבָּת׃}
}
{
	\eDtrA{
		and you shall remember that you were a slave in the land
		of Egypt and YHWH, your {\God}, brought you out from there
		with a strong hand and an outstretched arm. On account of
		this, YHWH, your {\God}, has commanded you to do the Sabbath
		day.
	}
}
{
}

\Verse{16}
{
	\hDtrA{כַּבֵּד אֶת אָבִיךָ וְאֶת אִמֶּךָ כַּאֲשֶׁר צִוְּךָ יְהֹוָה אֱלֹהֶיךָ לְמַעַן יַאֲרִיכֻן יָמֶיךָ וּלְמַעַן יִיטַב לָךְ עַל הָאֲדָמָה אֲשֶׁר יְהֹוָה אֱלֹהֶיךָ נֹתֵן לָךְ׃}
}
{
	\eDtrA{
		Honor your father and your mother, as YHWH, your {\God},
		commanded you, so that your days will be extended and so
		that it will be good for you on the land that YHWH, your
		{\God}, is giving you.
	}
}
{
%	\footnote{
%		Honor your father and your mother. Children ask, “What if your parents
%		tell you to break one of the other commandments?” The question reveals
%		an essential fact: all of the Ten Commandments are given to two people.
%		The commandment to the child to honor one’s parents is also a
%		commandment to the parent to be worthy of honor. The parent cannot be
%		utterly dishonorable, and cannot put the child in a position in which
%		the child cannot possibly honor the parent. Similarly, the command not
%		to steal is also a command not to deny food to the poor so that he or
%		she might be obliged to steal. Footnote 5:16. and so that it will be
%		good for you. These words do not appear in the Decalogue at Sinai. Moses
%		adds them here. See the comment on Deut Footnote 4:22. Again Moses
%		conveys the idea that the commandments are a path back to the good, the
%		condition that prevailed at the time of creation, prior to the event in
%		Eden.
%	}
}

\Verse{17}
{
	\hDtrA{לֹא תִּרְצָ͏ח׃ {ס} וְלֹא תִּנְאָ͏ף׃ {ס} וְלֹא תִּגְנֹב׃ {ס} וְלֹא תַעֲנֶה בְרֵעֲךָ עֵד שָׁוְא׃}
}
{
	\eDtrA{
		You shall not murder. And you shall not commit adultery.
		And you shall not steal. And you shall not testify against
		your neighbor as a false witness.
	}
}
{
%	\footnote{
%		adultery. This is perhaps the most readily broken of the commandments
%		(except for coveting?). It is the commandment that good people break.
%		The temptation is extraordinary. The potential for hurt is substantial.
%		The potential for leading to other behavior, especially lying, is
%		tremendous as well.
%	}
}

\Verse{18}
{
	\hDtrA{וְלֹא תַחְמֹד אֵשֶׁת רֵעֶךָ {ס} וְלֹא תִתְאַוֶּה בֵּית רֵעֶךָ שָׂדֵהוּ וְעַבְדּוֹ וַאֲמָתוֹ שׁוֹרוֹ וַחֲמֹרוֹ וְכֹל אֲשֶׁר לְרֵעֶךָ׃}
}
{
	\eDtrA{
		And you shall not covet your neighbor’s wife, and you
		shall not long for your neighbor’s house, his field, or
		his servant or his maid, his ox or his ass or anything
		that your neighbor has.
	}
}
{
}

\Verse{19}
{
	\hDtrA{אֶת הַדְּבָרִים הָאֵלֶּה דִּבֶּר יְהֹוָה אֶל כּל קְהַלְכֶם בָּהָר מִתּוֹךְ הָאֵשׁ הֶעָנָן וְהָעֲרָפֶל קוֹל גָּדוֹל וְלֹא יָסָף וַיִּכְתְּבֵם עַל שְׁנֵי לֻחֹת אֲבָנִים וַיִּתְּנֵם אֵלָי׃}
}
{
	\eDtrA{
		“YHWH spoke these words to your entire community at the
		mountain from inside the fire, the cloud, and the nimbus,
		a powerful voice. And He did not add. And He wrote them on
		two tablets of stones and gave them to me.
	}
}
{
%	\footnote{
%		He did not add. Again Moses reminds the people that the commandments as
%		stated are not to be added to (see the comment on Deut 4:2). For two
%		millennia Jews have struggled with the matter of additions to the law.
%		Rabbinites and Karaites, Orthodox and Reform and Conservative Jews, have
%		differed over the limits of the law and over where Jews have gone too
%		far in adding to it.
%	}
}

\Verse{20}
{
	\hDtrA{וַיְהִי כְּשׁמְעֲכֶם אֶת הַקּוֹל מִתּוֹךְ הַחֹשֶׁךְ וְהָהָר בֹּעֵר בָּאֵשׁ וַתִּקְרְבוּן אֵלַי כּל רָאשֵׁי שִׁבְטֵיכֶם וְזִקְנֵיכֶם׃}
}
{
	\eDtrA{
		“And it was when you heard the voice from inside the
		darkness, as the mountain was burning in fire, that you
		came forward to me, all the heads of your tribes and your
		elders,
	}
}
{
}

\Verse{21}
{
	\hDtrA{וַתֹּאמְרוּ הֵן הֶרְאָנוּ יְהֹוָה אֱלֹהֵינוּ אֶת כְּבֹדוֹ וְאֶת גּדְלוֹ וְאֶת קֹלוֹ שָׁמַעְנוּ מִתּוֹךְ הָאֵשׁ הַיּוֹם הַזֶּה רָאִינוּ כִּי יְדַבֵּר אֱלֹהִים אֶת הָאָדָם וָחָי׃}
}
{
	\eDtrA{
		and you said, ‘Here, YHWH, our {\God}, has shown us His glory
		and His greatness, and we’ve heard His voice from inside
		the fire. This day we’ve seen that {\God} may speak to a
		human and he lives.
	}
}
{
}

\Verse{22}
{
	\hDtrA{וְעַתָּה לָמָּה נָמוּת כִּי תֹאכְלֵנוּ הָאֵשׁ הַגְּדֹלָה הַזֹּאת אִם יֹסְפִים אֲנַחְנוּ לִשְׁמֹעַ אֶת קוֹל יְהֹוָה אֱלֹהֵינוּ עוֹד וָמָתְנוּ׃}
}
{
	\eDtrA{
		So now why should we die? Because this big fire will
		consume us. If we go on hearing the voice of YHWH, our
		{\God}, anymore, then we’ll die.
	}
}
{
}

\Verse{23}
{
	\hDtrA{כִּי מִי כל בָּשָׂר אֲשֶׁר שָׁמַע קוֹל אֱלֹהִים חַיִּים מְדַבֵּר מִתּוֹךְ הָאֵשׁ כָּמֹנוּ וַיֶּחִי׃}
}
{
	\eDtrA{
		Because who, of all flesh, is there who has heard the
		voice of the living {\God} speaking from inside fire as we
		have and lived?
	}
}
{
}

\Verse{24}
{
	\hDtrA{קְרַב אַתָּה וּשְׁמָע אֵת כּל אֲשֶׁר יֹאמַר יְהֹוָה אֱלֹהֵינוּ וְאַתְּ תְּדַבֵּר אֵלֵינוּ אֵת כּל אֲשֶׁר יְדַבֵּר יְהֹוָה אֱלֹהֵינוּ אֵלֶיךָ וְשָׁמַעְנוּ וְעָשִׂינוּ׃}
}
{
	\eDtrA{
		You go forward and listen to everything that YHWH, our
		{\God}, will say, and you’ll speak to us everything that
		YHWH, our {\God}, will speak to you, and we’ll listen, and
		we’ll do it.’
	}
}
{
}

\Verse{25}
{
	\hDtrA{וַיִּשְׁמַע יְהֹוָה אֶת קוֹל דִּבְרֵיכֶם בְּדַבֶּרְכֶם אֵלָי וַיֹּאמֶר יְהֹוָה אֵלַי שָׁמַעְתִּי אֶת קוֹל דִּבְרֵי הָעָם הַזֶּה אֲשֶׁר דִּבְּרוּ אֵלֶיךָ הֵיטִיבוּ כּל אֲשֶׁר דִּבֵּרוּ׃}
}
{
	\eDtrA{
		“And YHWH listened to your words’ voice when you were
		speaking to me, and YHWH said to me, ‘I’ve listened to the
		voice of this people’s words that they spoke to you.
		They’ve been good in everything that they’ve spoken.
	}
}
{
}

\Verse{26}
{
	\hDtrA{מִי יִתֵּן וְהָיָה לְבָבָם זֶה לָהֶם לְיִרְאָה אֹתִי וְלִשְׁמֹר אֶת כּל מִצְוֺתַי כּל הַיָּמִים לְמַעַן יִיטַב לָהֶם וְלִבְנֵיהֶם לְעֹלָם׃}
}
{
	\eDtrA{
		Who would make it so, that they would have such a heart,
		to fear me and observe all my commandments every day, so
		it would be good for them and for their children forever!
	}
}
{
%	\footnote{
%		Who would make it so. I asked in The Hidden Face of God: should God have
%		to ask who when the answer seems so manifestly to be that God Himself
%		could make it so? Even if “Who would make it so” was a known expression,
%		it still seems inapplicable for God. All the classic English
%		translations reword this verse, eliminating the phrase “who would make
%		it so,” presumably because it seems so incongruous. But we cannot
%		translate the concept away. It suggests that, in the Torah’s conception,
%		even though God makes humans as God sees fit, humans are free thereafter
%		to follow their hearts. God splits seas, makes animals talk, and makes
%		water come from rocks, but God does not change the human heart. This in
%		turn sheds light in retrospect on what God does to Pharaoh’s heart
%		during the plagues. English translators have often made it “God hardened
%		Pharaoh’s heart,” but the Hebrew rather means “God strengthened
%		Pharaoh’s heart.” God does not change the king’s heart, but only gives
%		the Pharaoh strength to follow his own resolve.
%	}
}

\Verse{27}
{
	\hDtrA{לֵךְ אֱמֹר לָהֶם שׁוּבוּ לָכֶם לְאהֳלֵיכֶם׃}
}
{
	\eDtrA{Go say to them, “Go back to your tents.”}
}
{
}

\Verse{28}
{
	\hDtrA{וְאַתָּה פֹּה עֲמֹד עִמָּדִי וַאֲדַבְּרָה אֵלֶיךָ אֵת כּל הַמִּצְוָה וְהַחֻקִּים וְהַמִּשְׁפָּטִים אֲשֶׁר תְּלַמְּדֵם וְעָשׂוּ בָאָרֶץ אֲשֶׁר אָנֹכִי נֹתֵן לָהֶם לְרִשְׁתָּהּ׃}
}
{
	\eDtrA{
		And you, stand here with me so I may speak to you all the
		commandment and the laws and the judgments that you shall
		teach them so they’ll do them in the land that I’m giving
		them to take possession of it.
	}
}
{
}

\Verse{29}
{
	\hDtrA{וּשְׁמַרְתֶּם לַעֲשׂוֹת כַּאֲשֶׁר צִוָּה יְהֹוָה אֱלֹהֵיכֶם אֶתְכֶם לֹא תָסֻרוּ יָמִין וּשְׂמֹאל׃}
}
{
	\eDtrA{
		And you shall be watchful to do as YHWH, your {\God}, has
		commanded you. You shall not turn right or left.
	}
}
{
%	\footnote{
%		not turn right or left. This expression complements the command not to
%		add to or subtract from the commandments (4:2; 5:19). The emphasis is on
%		focusing on the law as given. Later, King Josiah will be singled out
%		from all the kings of Israel and Judah because “he did not turn right or
%		left” (2 Kings 22:2).
%	}
}

\Verse{30}
{
	\hDtrA{בְּכל הַדֶּרֶךְ אֲשֶׁר צִוָּה יְהֹוָה אֱלֹהֵיכֶם אֶתְכֶם תֵּלֵכוּ לְמַעַן תִּחְיוּן וְטוֹב לָכֶם וְהַאֲרַכְתֶּם יָמִים בָּאָרֶץ אֲשֶׁר תִּירָשׁוּן׃}
}
{
	\eDtrA{
		You shall go in all the way that YHWH, your {\God}, has
		commanded you, so that you’ll live, and it will be good
		for you, and you’ll extend days in the land that you’ll
		possess.’
	}
}
{
}
